\documentclass{llncs}

\usepackage{graphicx}
\usepackage{xcolor}
\usepackage{listings}
\usepackage{mips}
\usepackage{amsmath}
\usepackage{cancel}

\newcommand{\crule}[3][black]{\textcolor{#1}{\rule{#2}{#3}}}

\newcommand{\erel}{\textsf{equiv}}
\newcommand{\merel}[1]{\textsf{equiv}_{#1}}
\newcommand{\pequiv}[2]{#1 \ \erel \ #2}
\newcommand{\pmequiv}[3]{#1 \ \merel{#3} \ #2}
\newcommand{\pnequiv}[2]{#1 \ \cancel{\erel} \ #2}
\newcommand{\pnmequiv}[3]{#1 \ \cancel{\merel{#3}} \ #2}

\definecolor{CommentGreen}{rgb}{0,.6,0}
\lstset{
 language=[mips]Assembler,
 escapechar=@, % include LaTeX code between `@' characters
 keepspaces,   % needed to preserve spacing with lstinline
 basicstyle=\footnotesize\ttfamily\bfseries,
 commentstyle=\color{CommentGreen},
 stringstyle=\color{cyan},
 showstringspaces=false,
 keywordstyle=[1]\color{blue},    % instructions
 keywordstyle=[2]\color{magenta}, % directives
 keywordstyle=[3]\color{red},     % registers
 }
 
\lstset{language=[mips]Assembler}

\title{Tamarin: Concolic Disequivalence for MIPS}
\author{Abel Nieto}
\institute{
University of Waterloo\\
\email{anietoro@uwaterloo.ca}
}

\begin{document}

\maketitle

\begin{abstract}
TODO
\end{abstract}

\section{Introduction}

We are staring at two opaque black boxes laying at our feet. Each box has a narrow slot through which we can place items in the box, but we cannot quite see what is inside. They look approximately like this:

\vspace{1mm}
\crule{1.5cm}{1.5cm}, \crule{1.5cm}{1.5cm}
\vspace{1mm}

We know each box contains an animal, but we do not know which specific animal is in each one. We would like to find out if both boxes contain the same species of animal. Our solution is simple: we take two carrots, and drop one in each box through the slots.

After a while, a chewing sound emerges from the boxes. We peer into them and, indeed, it looks like the carrots were successfully eaten. Triumphantly, we declare that the boxes contain the same species of animal. The truth is altogether different:

\vspace{1mm}
\fbox{\includegraphics[width=1.5cm]{camel}}, \fbox{\includegraphics[width=1.5cm]{beaver}}
\vspace{1mm}

The boxes are assembly programs. The animals are the functions those programs compute. The carrot is unit testing. The task was to determine whether the programs were equivalent. And we failed at it. In this paper, we show a technique that is better than the carrot.

Program equivalence. The program is the specification. The complications of assembly language.

\section{Program Equivalence for MIPS}

Let us set up the problem a bit more formally. Consider the set $P$ of MIPS-assembly programs that satisfy two restrictions: they take as inputs only the values of registers $\$1$ and $\$2$, and when they stop executing we define their output to be (exclusively) the value of $\$3$. Other side effects, such as printing values to the screen, or system calls, are disallowed.

We can now define a relation $\erel \subseteq P \times P$ of equivalent programs. Given $P_1, P_2 \in P$, we say that $\pequiv{P_1}{P_2}$ (read ``$P_1$ is equivalent to $P_2$'') if, for all inputs $\$1$ and $\$2$, one of the following holds:
\begin{itemize}
\item Both $P_1$ and $P_2$ fail during execution (for example, due to a divide-by-zero error).
\item $P_1$ and $P_2$ stop with the same output in $\$3$. 
\end{itemize} 

For example, the two programs in Figure \ref{equivprogs} are equivalent.
 
\begin{figure}
\begin{multicols}{2}
\begin{lstlisting}[escapeinside={(*}{*)}]
# P_1
add $3, $1, $2
\end{lstlisting}
\vfill\null
\columnbreak
\begin{lstlisting}
# P_2
add $4, $1, $1
lis $5
42
sw $4, 0, $5
add $3, $1, $2
\end{lstlisting}
\end{multicols}
\caption{$\pequiv{P_1}{P_2}$}
\label{equivprogs}
\end{figure}

Notice that $\pequiv{P_1}{P_2}$ even though $P_2$ modifies the contents of the memory and an additional register ($\$4$), because:
\begin{itemize}
\item Both $P_1$ and $P_2$ terminate without errors.
\item The value of $\$3$ will be the same when they do so.
\end{itemize}

Unfortunately, even though $\erel$ captures an already-simplified notion of equivalence\footnote{For example, $\erel$ has a very narrow notion of output that excludes side effects.}, a decision procedure for it does not exist, due to Rice's theorem.

To get decidability back, we define a new class of relations $\merel{S} \subseteq P \times P$. We say that $\pmequiv{P_1}{P_2}{S}$ (read ``$P_1$ is $S$-equivalent to $P_2$'') if, for all inputs, one of the following holds:
\begin{itemize}
\item Either $P_1$ or $P_2$ does not stop within $S$ steps (we can think of each CPU cycle as one step).
\item Both $P_1$ and $P_2$ fail.
\item Both $P_1$ and $P_2$ stop with the same output. 
\end{itemize}

The $\merel{S}$ relation captures the notion that we cannot tell $P_1$ and $P_2$ apart by running them for at most $S$ steps. Figure \ref{sequiv} shows an example of two programs that are $S$-equivalent for $S=10$, but not equivalent. This is the case because $P_2$ loops while the counter is less than $42$, so with 10 steps in our ``budget'' we will have to stop $P_2$ before the loop is over and we can observe the different result.

\begin{figure}
\begin{multicols}{2}
\begin{lstlisting}[escapeinside={(*}{*)}]
# P_1
add $3, $1, $2
\end{lstlisting}
\vfill\null
\columnbreak
\begin{lstlisting}
# P_2
  add $4, $0, 1  # counter
  add $5, $0, 42 # upper bound
loop:
  slt $6, $4, $5 
  beq $6, $0, end
  add $4, $4, 1
  beq $0, $0, loop
end:
  add $3, $1, $1  
\end{lstlisting}
\end{multicols}
\caption{$\pmequiv{P_1}{P_2}{10}$, but $\pnequiv{P_1}{P_2}$}.
\label{sequiv}
\end{figure}

Given a fixed $S$, the $\merel{S}$ relation is decidable because there is a finite number of inputs to try, and for each input we only need to run the programs a finite number of steps.

We already saw that equivalence not always implies $S$-equivalence. However, the converse always holds. The following lemma shows that $\merel{S}$ over-approximates $\erel$.

\begin{lemma}
\label{overapprox}
$\forall S, P_1, P_2$, $\pequiv{P_1}{P_2} \implies \pmequiv{P_1}{P_2}{S}$.
\end{lemma}

\begin{proof}
Let $\pequiv{P_1}{P_2}$. Then we have one of two cases:
\begin{itemize}
\item Either $P_1$ or $P_2$ (or both) do not stop within S steps. Then by definition $\pmequiv{P_1}{P_2}{S}$.
\item Both $P_1$ and $P_2$ stop within S steps. Then because they are equivalent, we know that they either fail with an error, or both stop with the same output. In either case, $\pmequiv{P_1}{P_2}{S}$.
\end{itemize}
\end{proof}

\begin{corollary}
\label{soundness}
$\pnmequiv{P_1}{P_2}{S} \implies \pnequiv{P_1}{P_2}$. 
\end{corollary}

\begin{proof}
This is just the contrapositive of Lemma \ref{overapprox}.
\end{proof}

Corollary \ref{soundness} can be used to argue the soundness (with respect to $\erel$) of any decision procedure that over-approximates $\merel{S}$. In the next section we will show one such over-approximation based on concolic execution.

\section{Concolic Disequivalence}


\section{Tamarin}

Overview.

\subsection{Trace Collection}

CPU instrumentation, PC concretization, error boxing, and fuel.

\subsection{Transformations}

Desugaring, simplification, trimming, and conversion to SSA.

\subsection{Query Representation}

Memory, jumps, arithmetic operators.

\subsection{Concolic Execution Redux}

Alternation. Compatibility. Soundness/Completeness. Efficiency.

\section{Evaluation}

\section{Related Work}

\section{Conclusions}

%\bibliographystyle{plain}
%\bibliography{refs}


\end{document}
